% year: תשפ"ו
% semester: א
% moed: ב
% number: 3
% points: 16
% categories: taylor
% official_solution: yes
% source: exams/מבחנים/תשפו/sol B.pdf

\begin{question}
(16 נק') חשבו בעזרת קירוב טיילור את $\sqrt[3]{11}$, כך שהשגיאה תהיה קטנה מ-$\frac{1}{10}$. נמקו.

שימו לב שהקירוב וחסם השגיאה צריכים להיות מספרים רציונליים.
\end{question}

\begin{hint}
פתחו את $f(x) = \sqrt[3]{x}$ סביב $a = 8$ (כי $\sqrt[3]{8} = 2$). התחילו מפולינום מסדר ראשון ובדקו אם שארית לגרנז' כבר קטנה מ-$\frac{1}{10}$.
\end{hint}

\begin{solution}
נגדיר $f(x) = \sqrt[3]{x} = x^{1/3}$ ונפתח סביב $a = 8$, הנקודה הקרובה ל-$11$ שבה השורש השלישי ידוע ($\sqrt[3]{8} = 2$). נחפש את $n$ המינימלי כך ש-$|R_n(11)| < \frac{1}{10}$.

\textbf{נגזרות} (עבור $x > 0$):
\[ f'(x) = \frac13 x^{-2/3}, \qquad f''(x) = -\frac29 x^{-5/3}, \]
ולכן $f(8) = 2$ ו-$f'(8) = \frac{1}{3\cdot 8^{2/3}} = \frac{1}{3\cdot 4} = \frac{1}{12}$.

\textbf{שארית לגרנז' עבור $n = 1$:} לפי משפט טיילור, קיימת $c$ בין $8$ ל-$11$ כך ש-
\[ R_1(11) = \frac{f''(c)}{2!}(11 - 8)^2 = \frac{-\frac29 c^{-5/3}}{2}\cdot 9 = -\frac{1}{c^{5/3}}. \]
מכיוון ש-$8 < c < 11$, מתקיים $c^{5/3} > 8^{5/3} = 2^5 = 32$, ולכן
\[ |R_1(11)| = \frac{1}{c^{5/3}} < \frac{1}{32} < \frac{1}{10}. \]
מכאן שפולינום טיילור מסדר ראשון מספיק ($n = 1$).

\textbf{הקירוב:}
\[ P_1(x) = f(8) + f'(8)(x - 8) = 2 + \frac{1}{12}(x - 8), \]
\[ \sqrt[3]{11} = f(11) \approx P_1(11) = 2 + \frac{3}{12} = 2\frac14 = \frac94. \]

\textbf{תשובה:} $\sqrt[3]{11} \approx \frac{9}{4} = 2.25$, עם שגיאה קטנה מ-$\frac{1}{32} < \frac{1}{10}$. (מכיוון ש-$R_1(11) < 0$, הקירוב גדול מהערך האמיתי: $\frac94 - \frac{1}{32} < \sqrt[3]{11} < \frac94$.)

(בדיקה: $\sqrt[3]{11} \approx 2.22398$, והשגיאה בפועל כ-$0.026 < \frac{1}{32} \approx 0.031$.)
\end{solution}
